% !PRINTED 251
\weberpage{251}
%<<CONT
gen\"ugen, und wenn man also $\gamma - \gamma_0 = \alpha\xi$ setzt (falls $\alpha$ von Null verschieden ist), so folgt $\delta - \delta_0 = \beta\xi$, also
\[ \gamma = \gamma_0 + \alpha\xi, \quad \delta = \delta_0 + \beta\xi, \]
und umgekehrt gen\"ugt auch jedes in dieser Form enthaltene Zahlenpaar $\gamma, \delta$. Dies gilt offenbar auch noch in dem ausgenommenen Falle $\alpha = 0$. Da wir nun $p^n$ verschiedene Zahlen f\"ur $\xi$ setzen k\"onnen, so ergiebt sich der Grad unserer Gruppe gleich
\weq{4}{p^n (p^{2n} - 1).}

Ist $p$ ungerade, so k\"onnen wir eine noch kleinere Gruppe ableiten:

Ist n\"amlich $p = 2$, so ist eine Zahl $\alpha$ von $-\alpha$ nicht verschieden, es sind also auch die Elemente
\weq{5}{\begin{pmatrix} \alpha, & \beta \\ \gamma, & \delta \end{pmatrix}, \quad \begin{pmatrix} -\alpha, & -\beta \\ -\gamma, & -\delta \end{pmatrix}}
mit einander identisch. Wenn aber $p$ ungerade ist, so sind die beiden Elemente (5) von einander verschieden, und das ganze System dieser Substitutionen zerf\"allt in Paare von der Form (5).

Wenn wir zwei solche Paare nehmen und componiren je ein Element des einen Paares mit je einem Element des anderen nach der Regel (2), so erhalten wir nur zwei verschiedene Elemente, die wieder ein solches Paar bilden. Diese Paare k\"onnen also selbst als Elemente einer neuen Gruppe vom Grade
\weq{6}{\dfrac{p^n (p^{2n} - 1)}{2}}
aufgefasst werden. Wir k\"onnen uns auch so ausdr\"ucken, dass zwei Elemente $A$, deren entsprechende Zahlen nur durch das Vorzeichen unterschieden sind, als nicht verschieden anzusehen sind.

Die so definirte Gruppe, deren Grad f\"ur $p = 2$ durch (4) und f\"ur ein ungerades $p$ durch (6) ausgedr\"uckt ist, wollen wir \gesperrt{die zum K\"orper $\mathfrak{C}$ geh\"orige Congruenzgruppe nennen und mit $E$ bezeichnen}.

F\"ur die genauere Untersuchung dieser Gruppe ist es von Wichtigkeit, ein System erzeugender Substitutionen zu kennen. Stellen wir nach \S.\,64, (2) die Zahlen von $\mathfrak{C}$ in der Form dar:
\weq{7}{\xi = x_0 + x_1 \varepsilon_1 + \cdots + x_{n-1} \varepsilon^{n-1},} % [sic] first power printed as subscript \varepsilon_1, not \varepsilon or \varepsilon^1
worin die $x_i$ rationale, nach dem Modul $p$ genommene Zahlen
%CONT>>
